% !TeX root = main.tex
\section{Background and Motivation}
\label{s:background}

In this section, we briefly describe the basics of volumetric video before describing the motivation for \sysname.

\subsection{Volumetric Video}

Volumetric video captures a three-dimensional representation of a continuously evolving scene (\eg a game, a live performance, a conference call).
%
Each viewer of a volumetric video has 6 degrees of freedom (6DoF)~---~he/she can view the video from any point in the captured space, in any direction.
%
For example, a viewer can watch a live performance of a string quartet as though sitting next to a violinist, while at the same time, another viewer can take the perspective of the cellist.
%
Typically, viewers watch volumetric videos on a mixed-reality headset.
%
They \emph{interact} with the video by moving around in a physical space; the
headset tracks the viewer's movement, and at any given moment, renders the
video from a perspective consistent with the viewer's current position and
orientation.
%
Thus, in our example of the musical performance, viewers can move to change perspective during the performance.
\ramesh{Consider adding a figure to show an RGB-D array surrounding a performance, and maybe show two different perspectives at the same time}

A volumetric video consists of a sequence of 3D frames, each of which represents a three-dimensional capture of a physical space.
%
There exist several three-dimensional frame representations; we focus on \emph{point clouds}~\cite{cite:vivo, cite:groot} due to their popularity, but overview alternate representations in \S\ref{s:related}.
%
Each point in a point cloud corresponds to a 3D location in physical space (usually the surface of an object) and is associated both with the coordinates of the location and color\footnote{Some sources of volumetric video, like LiDAR, do not produce color. We do not consider these.}
of the surface at that location.
%
Collectively, points in a point cloud capture the \emph{geometry} and \emph{color} of a scene in three dimensions. %\raj{Do we need to differentiate briefly with lidar point clouds with no color? Perhaps the angle that one is for downstream computation while color ptcl is for vieweing.}
 \rn{Other reasons to focus on point clouds? I added popularity but perhaps generation overheads also matter?}

In this paper, we consider volumetric videos obtained from an array of off-the-shelf RGB-D cameras encircling a scene. %\ramesh{If space permits, add a graphic for this as part of the earlier figure.}
%
Commodity RGB-D cameras, such as an Azure Kinect~\cite{cite:kinect} or an Intel RealSense~\cite{cite:realsense} camera, can capture both color (RGB) and depth (D) of points in a scene.
%
For example, the Kinect estimates depth using Time-of-flight~---~the time taken to receive the reflection of an emitted infrared beam.
%
By aligning (\cref{ss:view-generation}) RGB-D frames from multiple cameras, one can estimate the position and color of an arbitrary point in the scene, and thereby generate a point cloud.
%
Over the past few years, RGB-D cameras have become cheap enough that they are now a popular way to capture volumetric videos~\cite{xxx}.

% A volumetric video is a seqeunce of 3D captures represented explicitly as point clouds (ptcl)~\cite{cite:vivo, cite:groot} or meshes~\cite{cite:meshreduce}, implcitly represented using NeRFs~\cite{cite:nerf}. Out of these, point cloud have become quite popular, due to it's simplicity and flexibility. For every frame in a volumetric video, a point cloud is generated from RGB-D captures from multiple cameras (typically >$3$) synchronized together. Kinect synchronizes color and depth internally, while it uses external signal via 3.5-mm synchronization port, which the master camera uses to instruct the subordinate cameras to capture together. 

% \parab{Ptcl Generation.}
% Each synchronized color and depth frame from a camera is projected in 3D to form a point cloud using the camera's \textit{instrinsic parameters}. The point cloud is generated in the camera's local coordinate system by the following:
% \raj{Give the math behind point cloud generation.}
% A complete frame of volumetric video is generated by stitching the local point clouds from each camera in global coordinate system. Each of the camera's global position and rotation is estimated using extrinsic camera calibration methods~\cite{cite:zhang}~\cite{cite:icp}. Each point from the local point cloud is projected to the global coordinates by translation and rotation as obtained by the local camera's \textit{extrinsic parameters}. This is performed for all the points across all the cameras to form the stiched point cloud. A sequence of such sticted point clouds form a volumetric video.


\subsection{Truly Immersive Two-Way Live Streaming}
\label{s:immersive-two-way}

We consider the problem of streaming volumetric video between two locations (\textbf{A} and \textbf{B}) while permitting 6DoF viewing at both ends.
%
In contrast to prior work~\cite{starline,farfetchfusion} on live volumetric video streaming that transmits 3D representations of a single human being (their face or torso), we seek to stream an \textit{entire scene} from \textbf{A} that may consist of multiple participants and objects surrounding them (\textit{e.g.}, furniture, the floor, walls, etc.).
%
As users at \textbf{B} receive and view the video, they can move around the scene to change their perspective, permitting true immersion.
%
%\raj{6DoF allows users to have natural conversation as if they are same real-life scene. Should we differentiate early on that current systems like Starline and FarfetchFusion enjoy 3D conferencing but users are remain seated by design, thereby restricting 6DoF?}
%
These movements are captured in video transmitted to \textbf{A}, bringing a level of interactivity for human interaction that is comparable to physical presence.
%Further, receivers at \textbf{A} can see this movement in the video transmitted to them.
%
%This level of interactivity allows human interaction comparable to physical presence. 

To maximize the likelihood of Internet-scale adoption, we impose two requirements on a system supporting this capability:
\begin{description}
\item[Performance:] Live-streamed volumetric video must incur end-to-end latencies in the range of 200-400~ms~\cite{cite:acm-mm,salsify} and support 30 frames per second.
%
Today's 2-D video conferencing systems are designed for these latency targets because higher latencies have been shown to impair interactivity and user-perceived quality~\cite{xxx}.
\item[Bandwidth-adaptivity:] Live-streamed volumetric video must also adapt to changes in bandwidth availability, while minimally impacting quality.
%
In the absence of bandwidth-adaptivity, frames may not be delivered in time for rendering, resulting in \textit{stalls} that can mar user experience.
\end{description}

Prior research has not considered bandwidth-adaptive live streaming of 6DoF volumetric video (\cref{s:related}).
%
Most prior work has explored \emph{on-demand} 6DoF streaming of \emph{stored} volumetric video~\cite{xxx}.
%
This setting admits the possibility of significantly pre-processing the videos to ensure delivery quality (\eg as in Vivo~\cite{cite:vivo}, Veus~\cite{cite:veus}).
%
Work that has examined live streaming of volumetric videos (Holoportation~\cite{cite:holoportation}, LiveScan3D~\cite{cite:livescan3d}, Project Starline~\cite{cite:starline}, and MetaStream~\cite{cite:metastream}) has either constrained capture (\eg of faces or the torso, as described above), or is not bandwidth-adaptive, or both.
  
\subsection{Challenges}
\label{s:challenges}

High-quality volumetric video live streaming is challenging
because, relative to 2-D live streaming, stream quality is more susceptible to
stalls, end-to-end latency violations, or significant compression loss.
%
Each of these can degrade the visual quality of the reconstructed point cloud.
%
Objective point cloud quality metrics\footnote{These are the analogs of SSIM and PSNR, respectively, for 2D images.} such as PointSSIM~\cite{pointssim} or PointPSNR~\cite{p2psnr} can measure these degradations.
%
\cref{s:evaluation} describes these metrics in detail.
\rn{the para above felt a bit out of place to me, but I suspect I'm missing a tie.}


\parae{Large data volumes.}
%
Point clouds can have upwards of a million points~\cite{cite:8i-dataset, cite:vsense}.
%
Since each point is associated with a 3-D position and 3 color channels, the
total size of each frame can be larger than 15~MB.
%
At 30~fps, a volumetric video of these point clouds generates 3.6~Gbps of raw data.
%
This can stress both compute resources and network bandwidth.
%
In particular, without careful design, processing delays, such as those
required only for encoding and decoding point clouds, can result in stalls or violations of 
the end-to-end latency targets discussed above.
%
To combat this, as described earlier, prior work has considered constrained capture, resulting in smaller point cloud sizes (\cref{tab:sizes}).


In contrast, we consider transmitting point clouds that are an order of magnitude larger.
%
Average global broadband speeds are already nearly 100~Mbps and are expected to increase by 20\% annually~\cite{cisco}. Thus, we expect immersive two-way live streaming to be feasible in broadband scenarios in the near future, if it is not already.
\rn{I wonder if the bandwidth forecast should come in 2.2 (as a motivation to our focus on transmitting scenes)?}

\begin{table}[t]
  \begin{tabular}{|l|l|} \hline
    \textbf{Prior Work} & \textbf{Avg. Point Cloud Size (MB)} \\ \hline
  \end{tabular}
  \caption{Point cloud sizes used in prior work, compared to ours}
  \label{tab:sizes}
\end{table}

\parae{Inefficient point-cloud compression techniques.}
%
\antonio{This should be fleshed out. There are point cloud compression (PCC) methods, video-based (VPCC), and geometry-based (GPCC), both of which start from the point cloud. So, I think an argument we can make here is that our method is more efficient because the point cloud is only constructed at the decoder.  }
2-D video compression achieves significant compression efficiency using intra-frame and inter-frame compression.
%
In contrast, existing point cloud compression~\cite{draco,vpcc} does not achieve comparable levels of inter-frame compression, and can be slow.
%
Effective inter-frame point cloud compression is still the subject of active research~\cite{xxx}; it is a qualitatively harder problem than 2-D inter-frame compression because point clouds lack the spatial regularity of an array of pixels in 2-D images. \rn{could we mention results for this? both the pace of compression and efficacy?}

\parae{Lack of rate-adaptive codecs.}
%
%Unfortunately, point cloud compression codecs are not, to our knowledge, rate-adaptive.
%
2-D conferencing systems rely on rate-adaptive codecs that compress each frame to a target bandwidth estimated, for example, by the receiver~\cite{cite:salsify}.
%
This enables them to be bandwidth-adaptive.
%
\antonio{Actually, VPCC techniques are based on video codecs, so they can be rate-adaptive, but as pointed out above they require point cloud construction first and are pretty complex too. }
Unfortunately, to our knowledge, such rate-adaptive codecs have not been developed for point clouds.
%
Moreover, when they are developed, they will likely need to compress bandwidth to a low quality, given the inefficiency of inter-frame point cloud compression. 

\subsection{Approach}
\label{s:approach}

\rn{much of this seems redundant with Section 3.1 - perhaps merge into 3.1 and end section 2 with a short takeaway para?}

In this paper, we describe the design of \sysname, a bandwidth-adaptive live 6DoF volumetric video streaming system that satisfies the requirements listed above using two key ideas:
\begin{description}
\item[Using 2D video compression:] \sysname embeds the point cloud into a 2-D video stream, so it can leverage 2D video compression standards.
%
These, because they leverage inter-frame prediction, can be more efficient, and therefore higher quality, than point cloud compression under varying bandwidth conditions.
%
Moreover, 2D video codes are rate-adaptive, so this approach naturally supports bandwidth-adaptivity.
%
Even so, this is challenging for \sysname because 2-D video compression standards ensure high quality compression of \textit{color}.
%
To leverage 2-D video compression, \sysname must also encode \textit{geometry} (\ie depth) information in a 2-D image.
%
Volumetric video users can be sensitive to distortion in depth, so \sysname must compress depth information carefully.
\rn{this also implies that point cloud construction must happen at the receiver, right? perhaps that doesn't matter in this flow because we discuss 2-way communication (with the same compute restrictions on both sides)?}


\item[Frustum culling:] To reduce point cloud data volume, \sysname \textit{culls} points outside the \textit{frustum} (\cref{fig:architecture}, the 3-D equivalent of the \textit{viewport} in 2-D and 360-degree video~\cite{dragonfly}).
%
This is beneficial in two ways: (a) it can reduce compute latency if many points are culled, (b) relative to transmitting the full point cloud, the frustum can be transmitted at a higher quality at a given bandwidth.
%
To be able to cull, \sysname needs to predict the frustum accurately.
%
Frustum prediction is harder than viewport prediction; the former must predict motion along 6 degrees of freedom, while the latter considers only 3 dimensions (rotation in 360-degree video).
%
Prior work has used frustum and viewport (\cite{cite:flare, cite:dragonfly}) prediction, but only for on-demand streaming.
%
Fortunately, as we discuss later, in our setting, we can exploit the low end-to-end latency requirement to accurately predict frustums.
\rn{does the on-demand nature of prior work make this prediction easier somehow? is it just that it is more tolerant of errors since it has precomputed things so it gets back a bit more time to send more?}


%\rn{add text about fact that we can use short-term prediction (the opportunity/insight)}
\end{description}

% Today, video streaming has become commonplace in our lives, from entertainment, video-conference to remote operations. Though volumetric video streaming is more realistic and immersive compared to it's 2D counterpart, very few solutions streaming solutions exist both in academia and industry. Due to high data volume of streaming volumetric content, network bandwidth turns out to be the bottleneck. Each point in a colored point cloud is represented atleast by position in x, y and z-axis and 3 color channels as R, G, B and optional attributes such as faces, normals,~\etc. Position is represented in integer or floating point takes 12 bytes ($3 \times 4$ bytes) and color is represented as unsigned character takes 3 bytes ($3 \times 1$ byte), which is 15 bytes in total for a point. A ptcl frame having 1 million points~\cite{cite:8i-dataset} requires a raw size of 15~MB. Given standard streaming rate of 30~fps, the raw bandwidth required for streaming a volumetric video is 3.6~Gbps, which exceeds the available streaming bandwidth today.

% \parab{Live Streaming.}
% Volumetric video opens up opportunities for immersive applications such as immersive lectures, telepresence, remote surgery and live performances -- concerts, volumetric movies and sports. Immersive lectures in 6DoF allows remote students to experience live classroom with fellow classmates, walk up to the instructor to ask questions and present final projects in 3D, while the instructor (can also be remote) feels more engaged with students, demonstrate course material with immersive demos. But, these applications requires several steps such as scene capture, content generation, compression, transmission, decompression and rendering to be performed live. 

% Live streaming of volumetric content over the internet pose several requirements.

% \parae{Bandwidth Adaptation.}
% If available bandwidth is lower between sender and receiver than size of the video, the viewer experiences \textit{stalls}, where the current video frame freezes till next frame is received. On-demand video streaming services such as YouTube and Netflix use techniques such as receiver-side buffering and adapt to bitrate by switching between different pre-defined quality levels based on live receiver-side bandwidth estimation~\cite{cite:}. Live video conferencing services such as Zoom, Teams and Google Meet~\cite{cite:} cannot use either of the 2 approaches since buffering increases end-to-end latency, bitrate adaptation at pre-defined levels is computationally expensive for the sender. Instead today's conferencing systems uses low-latency rate controlled mode for the sender-side video encoder where encoder bitrate is changed dynamically based on the bandwidth estimate from the receiver~\cite{cite:salsify}. Similar approaches for bandwidth adaptation is more crucial in case of live volumetric video streaming systems due to large volume of data. There are several on-demand volumetric video streaming systems such Vivo~\cite{cite:vivo}, Veus~\cite{cite:veus} which propose several optimizations on the pre-recorded videos on the server to reduce the amount of transmitted data dynamically at the expense of computation. Few attempts on live volumetric video streaming such as Holoportation~\cite{cite:holoportation}, LiveScan3D~\cite{cite:livescan3d}, Project Starline~\cite{cite:starline} and MetaStream~\cite{cite:metastream} consider high available bitrate to transmit enire data or can adapt only to fixed set of quality levels \tbd{(Double check MetaStream's approach for adaptation)}.

% \raj{There are 2 types of live streaming: live streaming~\eg sports streaming and ultra-live streaming~\eg video conferencing. Often live streaming is implemented using chunk-based DASH with 2-10 sec of delay, while ultra-live streaming uses WebRTC to achieve ultra-low latency.}

% \parae{Low end-to-end latency.}
% End-to-end latency is the time between generation of a frame on the server and rendering of the same frame on the receiver. Current video conferencing services such as Zoom, Meet and Teams can achieve 200-300~ms of end-to-end latency~\cite{cite:acm-mm}. Similarly, a live volumetric streaming system should try to achieve latency closer to these production services. ~\raj{Should we mention why this latency should be small from human communication perspective?}

% \parae{High Quality.}
% Video streaming incur loss due to lossy compression, packets drops (partial frames) and incorrect captures resulting lower quality of experience (QoE). For 2D videos, the QoE is measured quantitatively using pixel-level root mean square error (RMSE), peak-signal-to-noise-ratio (PSNR) and structural similarity metric (SSIM)~\cite{cite:}. These metrics are correlated with qualitative evaluations by real users through carefully crafted user study. Similarly, 3D counterparts of these metrics are point-to-point RMSE (p2p-rmse), p2p-PSNR~\cite{cite:dong-liang} and PointSSIM~\cite{cite:pointssim} which have been prevalent in the signal processing community for a decade. A live volumetric streaming system must achieve high QoE both qualitatively and quantitatively.

% \parae{Easy deployable.}
% Today's video conferencing systems leverage investments in past decades from hardware to software to achieve high quality at low latency, low data rate while being robust against network dynamics. 2D video encoders such as H.265 (or HVEC), VP-9 have optimized implementations to leverage both CPU and GPU,~\eg Intel provides Intel VA-API~\cite{cite:vaapi}, Nvidia provides NVENC/NVDEC Codec SDK which are supported in open source video processing frameworks such as FFmpeg~\cite{cite:ffmpeg} and GStreamer~\cite{cite:gstreamer}. For live streaming, WebRTC has been a popular framework for low-latency congestion aware video streaming built over RTC and supported by most browsers, which can capture frames from sender's camera, encode, transmit and render at the receiver. Most popular videoconferencing services either use the above investments or have equivalent proprietary implementation inspired from these. An easily deployable live volumetric video streaming system can leverage these investments to make it work today. For example, ptcl encoders like Draco~\cite{cite:draco} and PCL~\cite{cite:pcl} either supports intra-frame compression only or have naive inter-frame compression, resulting in lower compression ratio. On contrary, any 2D video encoders use sophisticated motion compensation algorithms which can result in higher compression.

% \parab{Challenges.}
% Live volumetric video streaming faces several challenges as follows:
% \squishlist
%     \item At 30fps, the raw bandwidth requirement of volumetric video, where frames are represented as streams of ptcls, have a raw bandwidth requirement of 3.6 Gbps. Ptcl-based compression methods have \tbd{3x-4x} compression ratio, which still exceeds broadband bandwidth available to most household today~\cite{cite:broadband-study}. To this end, we need bandwidth adaptation techniques that works seamlessly with the compression technique we choose.
%     \item High compute requirements: \raj{Not sure if point cloud generation is a computationally expensive task. Rather should we say that live encoding and decoding is computationally expensive?}
% \squishend

% \parab{Strawman solutions.}
% In order to address some of the above challenges, we compile a set of strawman techniques from the current research in volumetric video streaming that can work live. 

% \parae{Compress Ptcl streams.} 
% The captured RGB-D frames from multiple cameras are stitched to form a ptcl. The sender can use current ptcl encoders such as Draco~\cite{cite:draco}, PCL~\cite{cite:pcl}, and V-PCC~\cite{cite:vpcc} to compress live streams of generated point clouds. The compressed binary stream can be transmitted to sender to the receiver over TCP. The receiver decodes to ptcl and renders on to a desktop, smartphone, or a head-mounted display (HMD). Ptcl compression techniques are still in their nascent stage with most techniques lacking motion compensation,  offer low compression ratio, high encoding latency and no hardware acceleration. A study of compressing full point clouds obtained from Panoptic Dataset~\cite{cite:panoptic} is shown in~\tabref{tab:ptcl-compression}.  

% \raj{Distinguish between lossy and loss-less compression, and whether existing lossy compression techniques allow for quality tradeoffs. What point are we trying to make?}

% \parae{Frustum Culling:}
% Previous works have observed that the data required by the viewer on the receiver-side is limited by the field-of-view (FoV) or frustum of the display device. The sender can cull the points outside the viewer's frustum, commonly known as \textit{frustum culling}, since sending or not sending the culled points do not improve the QoE. Though culling on the sender-side reduces data transmission rate, the viewer's frustum is known when the viewer consumes the frame. Thus, the sender needs to predict the frustum when the ptcl is generated, way ahead of it getting received and rendered. Frustum prediction schemes have been proposed for streaming 360-degree videos~\cite{cite:flare, cite:dragonfly}, but are still inaccurate. This becomes further challenging in 6DoF.

\begin{comment}

\parab{Outline}

********************************************
\raj{About 1-1.5 pages. For this section, as you write, if you can think of experiments/numbers to add, please note that down.}

\squishlist
\item Background on volumetric video:
    \squishlist
	\item How it is generated? Take example of multiple RGB-D cameras and show how depth and color are combined, and points are combined from multiple cameras. Example RGB cameras - kinect, realsense, etc. Sync cameras pointing inwards. Cost of sensors have gone down, thus we can deploy them in real settings.
	\item Video as a sequence of point cloud frames; can in theory be transmitted over the network and replayed at the other end (on a HMD, or stereoscopic display). What is point cloud? Data structure. How much each point takes, to bandwidth requirements at 30 fps. Why point cloud and not mesh. Easy to operate, simple and flexible repr - explicit representation than meshes or nerf (implicit).
        \item What is different about volumetric video? fundamentally interactive; clients can "move" to and watch from different parts of a scene
        \item Number of cameras in the datasets - N RGB or M RGBD cameras. 
    \squishend
	
\item Goal
\squishlist
    \item Live transmission of volumetric video generated from N RGB-D cameras.
    \squishlist
        \item Applications: immersive lectures, live performances, conferencing
    \squishend
	
    \item Requirements/constraints:
    \squishlist
        \item Adapt dynamically to available bandwidth
        \item Ensure relatively low end-to-end processing requirements (today's conferencing systems have latencies of 200-300ms)
        \item Preserve high quality
        \item When possible, leverage existing investments in live video streaming/conferencing (briefly describe how video conferencing works today).
        \squishlist
            \item Motion compensation is hard for live
            \item  Should this be a requirement or an insight?\
        \squishend
    \squishend
\squishend
	
\item Problem/Challenges:
\squishlist
    \item High bandwidth requirements for point clouds (get some ballpark numbers for this)
    \squishlist
        \item May need to adapt to a wide range of bandwidths (get rough numbers from various traces)
    \squishend
		
    \item High compute requirements (get some simple benchmarks for processing point clouds)
\squishend
	
\item Strawman approaches (not mutually exclusive):
\squishlist
    \item Compress point clouds: point cloud compression still an active topic of research, existing compression techniques do not give significant gains?
    \squishlist
        \item Distinguish between lossy and loss-less compression, and whether existing lossy compression techniques allow for quality tradeoffs
    \squishend
		
    \item View culling: just transmit part of the view; can be significantly expensive, and predicting the user's view is challenging, and mispredictions can impact quality of experience; viewport-adaptive approaches have been used for 360-video, but volumetric is different because users have 6 degrees of freedom
\squishend
\squishend

\raj{Related work: should have a para in the intro
(Microbenchmarks for justifying design choices: should we include them here?)}
\squishlist
    \item Explain the compute and network cost tradeoffs using some simple numbers/experiments - 
    \item Some observations to motivate our design.
\squishend
********************************************

\end{comment}
